
\documentclass{article}
\usepackage{colortbl}
\usepackage{makecell}
\usepackage{multirow}
\usepackage{supertabular}

\begin{document}

\newcounter{utterance}

\centering \large Interaction Transcript for game `cladder', experiment `full\_v1.5\_default', episode 8179 with Qwen3.5{-}9B.
\vspace{24pt}

{ \footnotesize  \setcounter{utterance}{1}
\setlength{\tabcolsep}{0pt}
\begin{supertabular}{c@{$\;$}|p{.15\linewidth}@{}p{.15\linewidth}p{.15\linewidth}p{.15\linewidth}p{.15\linewidth}p{.15\linewidth}}
    \# & \multicolumn{2}{c}{Player} && \multicolumn{2}{c}{Game Master} \\
    \hline

    \theutterance \stepcounter{utterance}  
    & & & \multicolumn{4}{p{0.6\linewidth}}{
        \cellcolor[rgb]{0.9,0.9,0.9}{
            \makecell[{{p{\linewidth}}}]{
                \texttt{\tiny{[P1$\langle$GM]}}
                \texttt{You are an expert in causal inference. The following question is not a typical commonsense query, but rather a meticulously designed question created by a professor specializing in causal inference, intended to assess the students' mastery of the course content.} \\
\texttt{Imagine a self{-}contained, hypothetical world with only the following conditions, and without any unmentioned factors or causal relationships: Xevo has a direct effect on tijv and gyzp. Tijw has a direct effect on tijv and gyzp. Tijv has a direct effect on gyzp. Tijw is unobserved. Method 1: We look directly at how xevo correlates with gyzp in general. Method 2: We look at this correlation case by case according to tijv. To understand how xevo affects gyzp, is it more correct to use the Method 1 than Method 2?} \\
\texttt{Start your answer with 'Yes' or 'No', followed by additional reasoning or evidence to support your explanation.} \\
            }
        }
    }
     \\ \\

    \theutterance \stepcounter{utterance}  
    & \multicolumn{4}{p{0.6\linewidth}}{
        \cellcolor[rgb]{0.561,0.737,0.561}{
            \makecell[{{p{\linewidth}}}]{
                \texttt{\tiny{[P1$\rangle$GM]}}
                \texttt{No} \\
\\ 
\texttt{**Reasoning:**} \\
\\ 
\texttt{In this hypothetical world, we must determine the causal effect of \$X\_\{evo\}\$ on \$G\_\{yzp\}\$. Let us first map the causal structure based on the provided conditions:} \\
\texttt{1.  \$X\_\{evo\} \textbackslash{}rightarrow T\_\{ijv\}\$} \\
\texttt{2.  \$X\_\{evo\} \textbackslash{}rightarrow G\_\{yzp\}\$} \\
\texttt{3.  \$T\_\{ijw\} \textbackslash{}rightarrow T\_\{ijv\}\$} \\
\texttt{4.  \$T\_\{ijw\} \textbackslash{}rightarrow G\_\{yzp\}\$} \\
\texttt{5.  \$T\_\{ijv\} \textbackslash{}rightarrow G\_\{yzp\}\$} \\
\texttt{6.  \$T\_\{ijw\}\$ is unobserved (a latent confounder).} \\
\\ 
\texttt{**Analysis of Method 1 (Total Effect):**} \\
\texttt{Method 1 looks at the correlation between \$X\_\{evo\}\$ and \$G\_\{yzp\}\$ without conditioning on any variables. In this specific DAG, \$X\_\{evo\}\$ affects \$G\_\{yzp\}\$ through two distinct pathways:} \\
\texttt{*   **Direct Path:** \$X\_\{evo\} \textbackslash{}rightarrow G\_\{yzp\}\$} \\
\texttt{*   **Indirect Path:** \$X\_\{evo\} \textbackslash{}rightarrow T\_\{ijv\} \textbackslash{}rightarrow G\_\{yzp\}\$} \\
\\ 
\texttt{The correlation observed in Method 1 captures the sum of these effects (the total causal effect). Since there are no confounders between \$X\_\{evo\}\$ and \$G\_\{yzp\}\$ (the only common cause of \$T\_\{ijv\}\$ and \$G\_\{yzp\}\$ is \$T\_\{ijw\}\$, which is not on the path from \$X\_\{evo\}\$ to \$G\_\{yzp\}\$), the marginal correlation correctly reflects the total impact of \$X\_\{evo\}\$ on \$G\_\{yzp\}\$.} \\
\\ 
\texttt{**Analysis of Method 2 (Conditional Effect):**} \\
\texttt{Method 2 conditions on \$T\_\{ijv\}\$ (looking at the correlation case by case according to \$T\_\{ijv\}\$). This corresponds to estimating the effect of \$X\_\{evo\}\$ on \$G\_\{yzp\}\$ *given* a specific value of \$T\_\{ijv\}\$.} \\
\texttt{*   By conditioning on \$T\_\{ijv\}\$, we block the indirect path (\$X\_\{evo\} \textbackslash{}rightarrow T\_\{ijv\} \textbackslash{}rightarrow G\_\{yzp\}\$).} \\
\texttt{*   The resulting correlation would only reflect the **direct effect** of \$X\_\{evo\}\$ on \$G\_\{yzp\}\$ (the path \$X\_\{evo\} \textbackslash{}rightarrow G\_\{yzp\}\$).} \\
\\ 
\texttt{**Conclusion:**} \\
\texttt{The question asks, "To understand how \$X\_\{evo\}\$ affects \$G\_\{yzp\}\$..." In causal inference, unless a specific direct effect is requested, "how X affects Y" generally refers to the **total causal effect**. Method 1 recovers the total effect, whereas Method 2 recovers only the direct effect by artificially blocking a valid causal mechanism (the mediation through} \\
            }
        }
    }
    & & \\ \\

    \theutterance \stepcounter{utterance}  
    & & & \multicolumn{2}{p{0.3\linewidth}}{
        \cellcolor[rgb]{0.9,0.9,0.9}{
            \makecell[{{p{\linewidth}}}]{
                \texttt{\tiny{[GM$|$GM]}}
                \texttt{No} \\
            }
        }
    }
    & & \\ \\

    \theutterance \stepcounter{utterance}  
    & & & \multicolumn{2}{p{0.3\linewidth}}{
        \cellcolor[rgb]{0.9,0.9,0.9}{
            \makecell[{{p{\linewidth}}}]{
                \texttt{\tiny{[GM$|$GM]}}
                \texttt{yes} \\
            }
        }
    }
    & & \\ \\

    \theutterance \stepcounter{utterance}  
    & & & \multicolumn{2}{p{0.3\linewidth}}{
        \cellcolor[rgb]{0.9,0.9,0.9}{
            \makecell[{{p{\linewidth}}}]{
                \texttt{\tiny{[GM$|$GM]}}
                \texttt{game\_result = LOSE} \\
            }
        }
    }
    & & \\ \\

\end{supertabular}
}

\end{document}
